% ============================================================================
% Cover Letter — Paper 4 (JoC / BMC)
% ============================================================================
\documentclass[11pt,a4paper]{letter}

\usepackage[margin=2.3cm]{geometry}
\usepackage{microtype}
\usepackage{hyperref}
\usepackage{enumitem}

\signature{Myke Vital Sao Temgoua \\ \small Corresponding Author}
\address{%
  University of Yaound\'{e} I \\
  Department of Physics, Faculty of Science \\
  P.O. Box 812, Yaound\'{e}, Cameroon \\
  \href{mailto:myke-vital.sao@facsciences-uy1.cm}{myke-vital.sao@facsciences-uy1.cm}
}

\begin{document}

\begin{letter}{%
  The Editors \\
  \textit{Journal of Cheminformatics} \\
  BioMed Central
}

\opening{Dear Editors,}

We are pleased to submit our manuscript, \textbf{``Pareto-guided Monte Carlo tree search explores multi-objective trade-offs missed by scalar antimalarial generation''}, for consideration as a Research Article in the \textit{Journal of Cheminformatics}.

The work addresses a recurring limitation of de novo molecular design: scalar reward aggregation that conceals trade-offs between potency, synthetic accessibility and drug-likeness. We introduce a Pareto-guided Monte Carlo tree search (MCTS) agent that maintains a non-dominated front across MPO, SYBA and SA scores, guided by a scaffold-aware fragment policy.

A four-method, 20-seed benchmark under a fixed oracle budget yields a deliberately honest ranking: random search achieves the highest mean scalar reward (0.665), followed by MCTS (0.659), a genetic algorithm (0.640) and deterministic greedy search (0.540). MCTS is therefore competitive on single-objective performance, yet is the only method that systematically returns a multi-objective Pareto front of non-dominated candidates (four solutions, hypervolume 1.237) spanning the potency--accessibility trade-off, including concave regions that scalar optimisation misses. Component ablations identify the chemistry-informed policy and the Pareto front itself as the dominant performance drivers.

We believe this transparent, methods-focused contribution aligns closely with the scope of the \textit{Journal of Cheminformatics}: rigorous benchmarking, open deposition of code and molecule sets, and honest reporting of both strengths and limitations of a multi-objective search strategy. All data and code will be deposited under the MIT licence.

This work is original, is not under consideration elsewhere, and all authors have approved the submission.

\closing{We thank you for your consideration.}

\ps{%
\textbf{Suggested reviewers:}\\
\begin{itemize}[leftmargin=*, itemsep=1pt]
  \item Dr.\ Jan H.\ Jensen (graph-based genetic algorithms / MCTS for chemical space)
  \item Prof.\ Gisbert Schneider (generative molecular design)
  \item Dr.\ Ola Engkvist (de novo design, multi-objective optimisation)
\end{itemize}
}

\end{letter}
\end{document}
